\documentclass{article}
\usepackage[T1]{fontenc}
\usepackage{lmodern}
\usepackage{graphicx}
\usepackage{amsmath,amssymb}
\usepackage[a4paper,margin=2cm]{geometry}
\usepackage[absolute,overlay]{textpos}
\setlength{\TPHorizModule}{1mm}
\setlength{\TPVertModule}{1mm}
\usepackage[indonesian]{babel}
\usepackage{hyperref}
\setlength{\emergencystretch}{2em}
% unit: Halaman 100

\begin{document}

% === Halaman 100 (llm) ===

\begin{align*}
K^{-1} &= \frac{1}{7} \begin{pmatrix} 3 & 10 \\ 15 & 9 \end{pmatrix} = 7^{-1} \begin{pmatrix} 9 & -10 \\ -15 & 3 \end{pmatrix} \\ &= 15 \begin{pmatrix} 9 & -10 \\ -15 & 3 \end{pmatrix} = 15 \begin{pmatrix} 9 & 16 \\ 11 & 3 \end{pmatrix} = \begin{pmatrix} 135 & 240 \\ 165 & 45 \end{pmatrix} \pmod{26} = \begin{pmatrix} 5 & 6 \\ 9 & 19 \end{pmatrix}
\end{align*}

\vspace{10mm}

\textbf{Keterangan} (ingat kembali teori bilangan di dalam Matematika Diskrit):
\begin{enumerate}[(i)]
\item $7^{-1} \pmod{26} \equiv 15$, karena $(7)(15) = 105 \pmod{26} = 1$
\item $-10 \equiv 16 \pmod{26}$
\item $-15 \equiv 11 \pmod{26}$ 
\end{enumerate}

\end{document}
